\section{Language-to-Rewards 的证据：从 Simulation 到 Real Robot}

机制成立后，需要检查三个问题：生成 reward 是否覆盖 diverse skills，是否优于更直接的 LLM interface，以及 simulator 中的 reward 与 controller 能否迁移到真实硬件。课程用 humanoid、quadruped、dexterous hand 和 real-robot manipulation 逐层回答。

\subsection{Prompt-to-skill：动作来自优化而非背诵}

Humanoid 与 manipulation 页面把 user prompt、生成 code 和最终 motion 放在一起。模型不需要在语言训练中见过每个关节轨迹；它利用“backflip”“open drawer”等语义写 state-based reward，optimizer 再搜索轨迹。因此同一 LLM 可通过不同 robot schema 生成不同 objective。

\lecturefigure{slide-47-humanoid-reward-code.jpg}{Humanoid skill：语言到 reward code，再到优化动作}{官方录像恢复幻灯片 V047；01:04:27}

\lecturefigure{slide-48-manipulation-reward-code.jpg}{Manipulation skill：open the drawer 的 reward 与动作}{官方录像恢复幻灯片 V048；01:04:39}

\begin{knowledgebox}{读 demo：至少检查四层}
第一层是 prompt 是否明确；第二层是 reward 是否真正编码任务而非画面 shortcut；第三层是 optimizer 是否稳定找到可行解；第四层是 execution 是否在扰动后保持成功。视频通常只展示第四层的成功片段，讲义必须把前面三层补回来。
\end{knowledgebox}

\subsection{Baseline comparison：reward-only 与 code-as-policy}

Benchmark 图比较 Language to Rewards、只由语言描述 reward 的简化方法，以及直接 code-as-policy 等 baseline。读图应看每个 task 的相对稳定性：reward interface 若在多种 morphology 上都优于 direct code，说明将 semantic objective 与 controller 分层具有价值；若少数任务差距小，则可能是 controller 本身已足够或 prompt 不需要复杂语义。

\lecturefigure{slide-49-language-to-rewards-benchmark.jpg}{Language to Rewards 与 reward-only / code-as-policy 比较}{官方录像恢复幻灯片 V049；01:05:06}

\begin{knowledgebox}{读图：比较的不是“谁更会写 Python”}
绿色、蓝色和黑色柱反映不同 interface 产生的最终 task performance。核心变量是 LLM 输出的抽象层级：reward specification 是否给优化器留下搜索空间，direct policy code 是否过早锁死动作。图支持 reward interface 在这些任务上更有效，但不说明所有任务都适合 MPC，也不衡量真实硬件延迟。
\end{knowledgebox}

\subsection{Sim-to-real：reward、trajectory 与 hardware 分层}

Sim-to-real 页面展示 reward translation、planned trajectory、simulator rollout 和 real execution。迁移成功依赖任务几何、robot calibration 与 controller robustness；如果 simulator dynamics 不准，MPC 可能产生在 simulation 中高分、真实世界不可执行的动作。Real-robot apple demo 因此是必要证据，但仍只是有限物体和平台上的验证。

\lecturefigure{slide-50-sim-to-real-transfer.jpg}{从 reward 与 planned trajectory 到 simulation / real execution}{官方录像恢复幻灯片 V050；01:05:36}

\lecturefigure{slide-51-real-robot-apple.jpg}{真实机器人执行“pick up the apple”}{官方录像恢复幻灯片 V051；01:05:57}

\begin{warningbox}{First use：sim-to-real gap}
\term{Sim-to-real gap} 是模拟与真实在视觉、动力学、接触、延迟、传感噪声和硬件磨损上的差异。常见缓解手段包括 domain randomization、system identification、robust control、online correction 和 real-data fine-tuning。一次成功 transfer 说明 pipeline 可行，不等于 gap 已消失。
\end{warningbox}

\subsection{LLM 作为 general pattern machine}

讲者最后把 LLM 抽象为 sequence transformation、sequence completion 与 sequence improvement。Language to Rewards 使用 transformation：instruction 变 reward code；interactive revision 使用 improvement：根据 feedback 修改 reward；RT-2 使用 completion：给定 image、instruction 与前缀 action，预测后续 token。这个抽象把两条路线放进同一计算框架。

\lecturefigure{slide-52-llm-general-pattern-machines.jpg}{LLM 的三类 pattern operation：转换、补全与改进}{官方录像恢复幻灯片 V052；01:07:12}

\lecturefigure{slide-53-sequence-improvement.jpg}{Sequence improvement：根据反馈迭代 reward 与动作}{官方录像恢复幻灯片 V053；01:08:51}

\teachervoice{课程把用户反馈放回 loop：如果生成的 motion 太快、姿态不自然或目标含糊，用户可以修改描述，LLM 更新 reward，controller 重新优化。这种交互性是 reward interface 的优势，但也意味着系统行为取决于 prompt、state schema、code validator 和 optimizer 的共同质量。}

\subsection{本章小结}

Language to Rewards 的证据链从 diverse simulated skills、baseline comparison 延伸到 sim-to-real 与 real robot。它说明 reward program 是一种有效语义接口，却仍依赖 controller、dynamics model、validation 和硬件安全。两条路线的差异将在下一章统一比较。

